\chapter{Adrenalectomy}

\section*{Introduction \& Clinical Indications}
Adrenalectomy is a surgical procedure that involves the removal of one or both adrenal glands, which are small, triangular-shaped endocrine glands located atop each kidney. These glands play a crucial role in the body's hormonal balance by producing hormones such as cortisol, aldosterone, adrenaline (epinephrine), and noradrenaline (norepinephrine). The surgery is primarily indicated for the treatment of adrenal tumors that may be functioning (producing excess hormones), suspected or confirmed malignant tumors, or large incidentalomas that pose a risk for malignancy or cause significant mass effect. The procedure can be performed using traditional open techniques or, more commonly today, through minimally invasive laparoscopic approaches, which offer advantages such as reduced postoperative pain, shorter recovery times, and improved cosmetic outcomes.

\begin{itemize}
  \item Adrenal gland resection
  \item Laparoscopic adrenalectomy
  \item Open adrenalectomy
  \item Unilateral adrenalectomy
  \item Bilateral adrenalectomy
  \item Adrenal tumor excision
\end{itemize}

The fundamental principle of adrenalectomy is the complete and safe removal of the affected adrenal gland or tumor, aiming to eliminate the source of pathological hormone overproduction or to resect a malignant or potentially malignant mass. For functioning tumors such as pheochromocytomas, aldosteronomas, or cortisol-producing adenomas, the rationale is to cure the associated endocrine syndrome, thereby resolving symptoms like hypertension, hypokalemia, or Cushingoid features. This involves meticulous dissection to isolate the gland, ligate its vascular supply, and extract it intact.

In cases of adrenocortical carcinoma or metastatic disease to the adrenal gland, the primary goal is oncological: achieving a complete R0 resection (no residual microscopic disease at the margins) to maximize the chance of cure or local control. This often necessitates a wider margin of resection and careful evaluation for local invasion or lymph node involvement. For large, non-functioning incidentalomas, the rationale for removal is to prevent potential future malignancy or to alleviate symptoms from mass effect, guided by imaging characteristics and growth patterns. The technical goal is always to minimize morbidity while achieving the therapeutic objective, often favoring minimally invasive approaches when appropriate.

The history of adrenalectomy is closely linked to the understanding of adrenal physiology and pathology. The first successful adrenalectomy for a functioning tumor (pheochromocytoma) was performed by C.H. Mayo in 1927 at the Mayo Clinic, following earlier attempts and increasing knowledge of adrenal pathology. Early operations were exclusively open procedures, requiring large abdominal or flank incisions, which were associated with significant morbidity and prolonged recovery. The management of pheochromocytoma, in particular, was fraught with high mortality due to intraoperative hypertensive crises until the advent of effective preoperative alpha-adrenergic blockade in the 1950s.

A major paradigm shift occurred in 1992 when Michel Gagner performed the first laparoscopic adrenalectomy. This minimally invasive approach rapidly gained acceptance due to its advantages, including reduced postoperative pain, shorter hospital stays, faster recovery, and improved cosmetic outcomes. Today, laparoscopic adrenalectomy, including the transabdominal and retroperitoneal approaches, is the standard of care for most benign adrenal tumors and many smaller malignant lesions. Open adrenalectomy is now reserved for very large tumors, those with extensive local invasion, or suspected adrenocortical carcinomas where complete resection requires a wider exposure.

Adrenalectomy is indicated for a variety of adrenal pathologies, primarily those causing hormonal excess or malignancy. Specific indications include:

\begin{itemize}
  \item Functioning adrenal adenomas causing primary aldosteronism (Conn's syndrome), characterized by hypertension and hypokalemia, often diagnosed via aldosterone-to-renin ratio and confirmed by adrenal vein sampling.
  \item Cortisol-producing adenomas leading to Cushing's syndrome, diagnosed by elevated cortisol levels and confirmed with dexamethasone suppression tests.
  \item Pheochromocytoma, a catecholamine-secreting tumor of the adrenal medulla, causing paroxysmal or sustained hypertension, palpitations, and headaches, diagnosed by elevated plasma or urine metanephrines.
  \item Adrenocortical carcinoma (ACC), a rare but aggressive malignancy, often presenting as a large, rapidly growing mass, diagnosed by imaging and confirmed by pathology.
  \item Metastatic disease to the adrenal gland from other primary cancers (e.g., lung, kidney, melanoma) in selected patients with oligometastatic disease.
  \item Large adrenal incidentalomas (typically >4-6 cm) that are non-functional but carry a risk of malignancy or demonstrate significant growth on serial imaging.
  \item Adrenal masses with suspicious imaging characteristics (e.g., high Hounsfield units on CT, irregular borders, rapid enhancement/washout) concerning for malignancy, regardless of size.
  \item Bilateral adrenal hyperplasia causing Cushing's syndrome or primary aldosteronism refractory to medical management, sometimes requiring bilateral adrenalectomy.
\end{itemize}

Adrenalectomy holds a primary and definitive role in the management of most localized, functioning adrenal tumors and resectable adrenocortical carcinomas. For benign, hormone-producing lesions such as aldosteronomas, cortisol-producing adenomas, and pheochromocytomas, surgical removal of the affected gland is considered the primary curative treatment, often leading to complete resolution of the endocrine syndrome. In these cases, medical management is typically used for preoperative optimization or as a temporizing measure, but surgery offers the best chance for a cure.

For localized adrenocortical carcinoma, complete surgical resection is the cornerstone of treatment and offers the only potential for long-term survival, making it a first-line intervention. In contrast, for bilateral adrenal pathologies, such as bilateral adrenal hyperplasia causing Cushing's syndrome, bilateral adrenalectomy may be considered a secondary or salvage procedure after medical therapies have failed or are poorly tolerated. Similarly, for metastatic disease to the adrenal gland, adrenalectomy is often an adjunct therapy, performed in conjunction with systemic treatments or as part of a multidisciplinary approach for selected patients with limited metastatic burden.

\section*{Relevant Surgical Anatomy}
The adrenal glands are paired retroperitoneal organs, each weighing approximately 4-5 grams, located superior and medial to the upper pole of the kidneys, typically at the level of the T11-L2 vertebrae. Each gland consists of an outer cortex and an inner medulla, which have distinct embryonic origins and hormonal functions. The right adrenal gland is pyramidal in shape, situated posterior to the inferior vena cava (IVC) and medial to the liver, often partially covered by the liver. Its short central vein drains directly into the IVC. The left adrenal gland is crescent-shaped, positioned medial to the spleen, superior to the pancreas, and lateral to the aorta, with its central vein draining into the left renal vein.

The arterial supply to each gland is rich, arising from three main sources: the superior adrenal arteries (branches of the inferior phrenic artery), the middle adrenal artery (a direct branch of the aorta), and the inferior adrenal artery (a branch of the renal artery). Venous drainage typically occurs via a single central vein from each gland. Lymphatic drainage follows the arterial supply, primarily to the para-aortic lymph nodes. The glands are innervated by preganglionic sympathetic fibers originating from the splanchnic nerves, particularly to the adrenal medulla, which functions as a modified sympathetic ganglion. Surgical landmarks include the superior pole of the kidney, the IVC on the right, the aorta on the left, and the crus of the diaphragm posteriorly. Careful dissection is required to avoid injury to these adjacent vital structures.

\section*{Preoperative Workup \& Preparation}
Thorough preoperative workup is essential to optimize patient safety and surgical outcomes. This includes a comprehensive medical and cardiac evaluation to assess fitness for general anesthesia and surgery, especially given the potential for significant hemodynamic shifts. Routine laboratory tests include a complete blood count, electrolyte panel, renal and liver function tests, and coagulation studies. Imaging, such as CT or MRI of the abdomen, is crucial to precisely localize the adrenal mass, characterize its features (size, density, enhancement pattern), and assess for local invasion or metastatic spread.

Specific endocrine preparation is paramount for functioning tumors:

\begin{itemize}
  \item **Pheochromocytoma:** Mandatory alpha-adrenergic blockade (e.g., phenoxybenzamine) for 10-14 days preoperatively, followed by beta-blockade if needed for tachycardia, to prevent a life-threatening hypertensive crisis during tumor manipulation.
  \item **Cushing's Syndrome:** Optimization of blood pressure, blood glucose, and electrolyte imbalances. Patients may require preoperative steroid tapering or blockade (e.g., metyrapone, ketoconazole) to reduce cortisol levels.
  \item **Primary Aldosteronism:** Control of hypertension and correction of hypokalemia with mineralocorticoid receptor antagonists (e.g., spironolactone) or potassium supplementation.
\end{itemize}

Patients are instructed to be NPO (nil per os) for at least 6-8 hours prior to surgery. Medications, particularly anticoagulants or antiplatelet agents, are adjusted or held according to established protocols. Prophylactic broad-spectrum antibiotics are administered intravenously within 60 minutes of incision. Informed consent is obtained, detailing the procedure, potential risks, benefits, and the possibility of requiring lifelong steroid replacement, especially if bilateral adrenalectomy is performed or if the contralateral gland is suppressed.

Adrenalectomy, while generally safe, has specific contraindications and requires careful precautions.

**Absolute Contraindications:**

\begin{itemize}
  \item Uncontrolled coagulopathy or severe bleeding disorder that cannot be corrected preoperatively.
  \item Severe, uncompensated cardiovascular or pulmonary disease that renders the patient an unacceptable anesthetic risk.
  \item Widespread metastatic disease where adrenalectomy would not offer significant palliative benefit or improve survival.
  \item Patient refusal or inability to provide informed consent.
\end{itemize}

**Relative Contraindications \& Precautions:**

\begin{itemize}
  \item Very large adrenal tumors (>10 cm) or those with extensive local invasion, which may necessitate an open approach rather than laparoscopic.
  \item Significant comorbidities (e.g., severe obesity, prior extensive abdominal surgery) that increase the technical difficulty or risk of a laparoscopic approach, potentially favoring an open conversion.
  \item Pregnancy, where surgery is generally deferred until after delivery unless there is an urgent indication (e.g., symptomatic pheochromocytoma).
  \item For pheochromocytoma, inadequate preoperative alpha-blockade is a critical precaution; proceeding without it significantly increases the risk of intraoperative hypertensive crisis and mortality.
  \item Patients undergoing bilateral adrenalectomy or those with a suppressed contralateral gland require meticulous planning for perioperative and lifelong steroid replacement to prevent adrenal insufficiency.
\end{itemize}

\section*{Surgical Technique \& Procedure}
Adrenalectomy is typically performed under general anesthesia. Patient positioning varies by approach: for a transabdominal laparoscopic approach, the patient is usually placed in a lateral decubitus position with the affected side up, allowing gravity to retract the abdominal contents. For a posterior retroperitoneal laparoscopic approach, the patient is often prone or in a modified lateral position.

The general steps for a laparoscopic adrenalectomy include:

\begin{itemize}
  \item **Anesthesia \& Positioning:** General endotracheal anesthesia is induced. The patient is positioned as described, and the surgical field is prepped and draped.
  \item **Access \& Insufflation:** Small incisions (typically 3-4) are made for trocar placement. For transabdominal, the abdomen is insufflated with carbon dioxide to create pneumoperitoneum. For retroperitoneal, a working space is created without entering the peritoneal cavity.
  \item **Exploration \& Dissection:** A laparoscope is inserted to visualize the abdominal cavity or retroperitoneal space. The adrenal gland is identified, typically superior to the kidney.
  \item **Mobilization:** The gland is carefully mobilized from surrounding structures (e.g., liver on the right, spleen/pancreas on the left, kidney inferiorly, diaphragm posteriorly).
  \item **Vascular Control:** The adrenal arteries and the central adrenal vein are identified, carefully dissected, clipped, and divided. This is a critical step to prevent hemorrhage and, in pheochromocytoma, to minimize catecholamine release.
  \item **Gland Removal:** Once completely free, the adrenal gland is placed into an endoscopic retrieval bag.
  \item **Specimen Extraction:** The bag and specimen are removed through one of the larger port sites, which may require slight enlargement.
  \item **Closure:** The port sites are inspected for bleeding. Fascial defects larger than 10 mm are closed to prevent hernia formation. Skin incisions are closed with sutures or staples. Drains are rarely used unless there is significant concern for fluid collection or hemorrhage.
\end{itemize}

For open adrenalectomy, a subcostal, flank, or midline incision is used, providing direct access for larger tumors or complex resections. The duration of the procedure typically ranges from 1 to 3 hours, depending on the complexity and approach.

\section*{Complications \& Risks}
Adrenalectomy carries both general surgical risks and specific procedure-related complications.

**General Surgical Complications:**

\begin{itemize}
  \item Bleeding, potentially requiring blood transfusion or reoperation.
  \item Infection, including wound infection, intra-abdominal abscess, or pneumonia.
  \item Complications related to general anesthesia (e.g., allergic reactions, cardiac events, respiratory depression).
  \item Deep vein thrombosis (DVT) and pulmonary embolism (PE).
  \item Incisional hernia at port sites (laparoscopic) or wound dehiscence (open).
\end{itemize}

**Procedure-Specific Complications:**

\begin{itemize}
  \item **Hypertensive crisis or severe hemodynamic instability:** Particularly during pheochromocytoma resection, due to catecholamine release upon tumor manipulation, or hypotension post-resection.
  \item **Injury to adjacent organs:**
    \begin{itemize}
      \item Right side: Liver, inferior vena cava (IVC), duodenum, kidney.
      \item Left side: Spleen (potentially requiring splenectomy), pancreas (risk of pancreatitis or fistula), kidney, stomach.
    \end{itemize}
  \item **Adrenal insufficiency/Addisonian crisis:** If both glands are removed, or if the remaining gland is suppressed and steroid replacement is inadequate. This is a life-threatening emergency.
  \item **Vascular injury:** To the renal vessels, aorta, or IVC, which can lead to significant hemorrhage.
  \item **Nerve injury:** Rarely, injury to intercostal nerves causing chronic pain, or phrenic nerve injury causing diaphragmatic dysfunction.
  \item **Pancreatitis:** More common with left adrenalectomy due to proximity of the pancreatic tail.
  \item **Recurrence of disease:** Especially for pheochromocytoma (up to 10-15\% for sporadic cases, higher for genetic syndromes) or adrenocortical carcinoma.
\end{itemize}

\section*{Postoperative Care \& Recovery}
The postoperative recovery timeline for adrenalectomy varies depending on the surgical approach (laparoscopic vs. open) and the patient's overall health. Immediately after surgery, the patient is transferred to the Post-Anesthesia Care Unit (PACU) for close monitoring of vital signs, pain levels, and fluid balance. For pheochromocytoma patients, blood pressure is meticulously monitored for hypotension following tumor removal.

Within the first 24-48 hours, patients undergoing laparoscopic adrenalectomy are typically encouraged to ambulate early, which aids in preventing DVT and promoting bowel function. Oral intake usually begins with clear liquids and progresses to a regular diet as tolerated. Pain is managed with oral analgesics, often transitioning from intravenous medications. Most laparoscopic patients are discharged within 1 to 3 days. During the first 1-2 weeks, patients are advised to avoid heavy lifting and strenuous activities. Full recovery, including return to normal activities and work, is generally expected within 2-4 weeks for laparoscopic surgery, while open adrenalectomy may require 6-8 weeks or longer due to larger incisions and greater tissue disruption. Long-term recovery involves ongoing monitoring of hormone levels and, for some, lifelong hormone replacement therapy.

During adrenalectomy, the surgeon meticulously documents intraoperative findings, including the size, location, and gross appearance of the adrenal mass, its relationship to surrounding structures, and any evidence of local invasion or lymphadenopathy. This information is crucial for staging and guiding further management. The resected adrenal gland and tumor are then sent for pathological examination.

The pathology report provides a definitive diagnosis, characterizing the tumor as benign (e.g., adenoma, pheochromocytoma, myelolipoma) or malignant (e.g., adrenocortical carcinoma, metastatic carcinoma). For benign lesions, the report confirms the tumor type and often the presence of hormonal activity. For malignant tumors, the pathologist assesses tumor size, histological subtype, mitotic rate, nuclear atypia, presence of necrosis, capsular invasion, vascular invasion, and surgical margin status (positive or negative). These pathological features are critical for determining the prognosis and guiding adjuvant therapy. The correlation between intraoperative findings and the final pathology report is essential for comprehensive patient care and accurate disease staging.

A normal and successful postoperative course following adrenalectomy is characterized by a smooth recovery with resolution of the underlying pathology. Patients typically experience a gradual decrease in incisional pain, which is well-controlled with oral analgesics. For functioning tumors, the expected course includes the normalization of hormone levels, leading to resolution of symptoms such as hypertension, hypokalemia, or Cushingoid features. Blood pressure and blood glucose levels should stabilize, often requiring adjustment of antihypertensive or antidiabetic medications.

Patients are expected to resume normal bowel function within a few days and gradually increase their activity levels. Wound healing progresses without complications, and any sutures or staples are typically removed at a follow-up visit. For patients requiring steroid replacement, the expected course involves establishing a stable and appropriate hormone regimen. The overall goal is a return to baseline health or an improved quality of life due to the resolution of the adrenal-related condition, with minimal long-term morbidity.

Postoperative care for adrenalectomy patients involves several key components to ensure optimal recovery and long-term management.

\begin{itemize}
  \item **Wound Care:** Instructions for keeping incision sites clean and dry. Suture or staple removal typically occurs 1-2 weeks post-surgery.
  \item **Pain Management:** Continued use of oral analgesics as needed, with gradual tapering.
  \item **Activity Restrictions:** Avoidance of heavy lifting and strenuous exercise for several weeks to prevent incisional hernia or wound complications.
  \item **Endocrine Follow-up:** Regular monitoring of hormone levels (e.g., cortisol, aldosterone, metanephrines) by an endocrinologist to confirm resolution of hormonal excess and to manage any necessary hormone replacement therapy.
  \item **Steroid Replacement:** For patients who have undergone bilateral adrenalectomy or whose contralateral gland is suppressed, lifelong glucocorticoid and mineralocorticoid replacement is initiated and carefully titrated. Patients are educated on "sick day rules" for steroid dose adjustment.
  \item **Imaging Surveillance:** For patients with adrenocortical carcinoma or pheochromocytoma, regular imaging (CT or MRI) and biochemical surveillance are crucial to monitor for recurrence or metastatic disease.
  \item **Warning Signs:** Patients are educated on symptoms requiring immediate medical attention, such as fever, worsening pain, redness or drainage from incision sites, persistent nausea/vomiting, severe fatigue, or symptoms of adrenal insufficiency (e.g., profound weakness, dizziness, hypotension).
\end{itemize}

Regular follow-up visits with the surgeon and endocrinologist are scheduled to monitor recovery, review pathology results, and adjust long-term management plans.

\section*{Outcomes \& Prognosis}
Adrenal masses are increasingly identified incidentally on abdominal imaging (adrenal incidentalomas), with a prevalence rising with age, affecting up to 10\% of individuals over 70 years old. The vast majority (70-80\%) of these incidentalomas are benign, non-functional cortical adenomas. However, a significant minority are functional or malignant. Primary aldosteronism, caused by an aldosteronoma or bilateral hyperplasia, is now recognized as a common cause of secondary hypertension, affecting 5-10\% of hypertensive patients and up to 20\% of those with resistant hypertension. Pheochromocytomas are rare, with an estimated incidence of 0.8 per 100,000 person-years, but are critical to diagnose due to their life-threatening potential. Adrenocortical carcinoma (ACC) is exceedingly rare, with an incidence of 0.7-2 cases per million population per year, but carries a poor prognosis. The overall mortality and morbidity associated with adrenalectomy are low, particularly with laparoscopic approaches, reflecting advancements in surgical technique and perioperative management.

The outcomes and prognosis following adrenalectomy are generally excellent for benign, functioning tumors. For primary aldosteronism due to an aldosteronoma, unilateral adrenalectomy leads to biochemical cure in nearly 100\% of patients and significant improvement or cure of hypertension in 50-70\%, as demonstrated in studies like the Primary Aldosteronism Surgical Outcome (PASO) study. Patients with Cushing's syndrome from an adrenal adenoma typically achieve biochemical remission and improvement in associated comorbidities (e.g., diabetes, osteoporosis) in over 90\% of cases. For pheochromocytoma, surgical resection is curative in 90-95\% of sporadic cases, though lifelong surveillance for recurrence is recommended, especially in patients with genetic syndromes (e.g., MEN2, VHL, NF1).

Adrenocortical carcinoma (ACC) has a more guarded prognosis, with 5-year survival rates ranging from 15-60\% depending on stage. Complete surgical resection (R0) is the most important prognostic factor for ACC, offering the best chance for long-term survival. Recurrence rates for ACC remain high, even after complete resection, necessitating aggressive surveillance and often adjuvant therapies. Overall, the shift to minimally invasive techniques has significantly improved short-term outcomes, reducing hospital stay and recovery time, while maintaining excellent long-term functional and survival outcomes for appropriate indications.

\section*{References}
\begin{enumerate}
  \item MedlinePlus. Adrenalectomy. U.S. National Library of Medicine; 2024.
  \item Sabiston DC, Townsend CM. Sabiston Textbook of Surgery: The Biological Basis of Modern Surgical Practice. 21st ed. Elsevier; 2021.
  \item Zeiger MA, Thompson GB, Duh QY, et al. American Association of Endocrine Surgeons Guidelines for the Management of Adrenal Incidentalomas: AACE/AAES Adrenal Incidentaloma Guidelines Task Force. Surgery. 2009;146(6):1188-1203.
  \item Fassnacht M, Arlt W, Bancos I, et al. Management of Adrenal Incidentalomas: European Society of Endocrinology Clinical Practice Guideline. Eur J Endocrinol. 2016;175(2):G1-G34.
\end{enumerate}

\clearpage
